% Section 1: Introduction
\section{Introduction}
\label{sec:introduction}

A parallel system of $m$ components functions as long as at least one
component survives; its lifetime is $T = \max(T_1, \ldots, T_m)$.  Such
systems arise in redundant engineering designs---backup power supplies,
server clusters, communication networks---where each component provides
an independent path to system functionality.  A basic engineering
question is: given observed system lifetimes, what can be learned about
the individual component lifetime distributions?

This question arises in practice whenever component-level testing is
impractical or impossible: satellites with redundant subsystems
observable only at system level, wind turbine gearboxes where
individual bearing instrumentation is costly, nuclear safety systems
tested by activating the entire redundant chain, and data center
servers whose component failures are recorded only at the machine
level.  In each case, system failure times are the primary data source,
and component-level inference requires exploiting the system structure.

For the dual topology---a series system with
$T = \min(T_1, \ldots, T_m)$---this question has generated a vast
literature under the heading of \emph{competing risks} and
\emph{masked cause of failure}.  The foundational result of
\citet{Tsiatis1975} establishes that component distributions are
\emph{not} identifiable from system data alone; additional
information about the cause of failure is needed.  Estimation under
partial cause information (masking) has been studied extensively
\citep{Miyakawa1984, UsherHodgson1988, LinUsherGuess1993,
GuessUsherHodgson1991, KunduBasu2000, CraiuReiser2006}.

The parallel case has received far less attention, despite sitting on
the opposite side of an identifiability divide.  \citet{BasuGhosh1980}
showed that the distribution of $\max(T_1, \ldots, T_m)$ uniquely
determines the component distributions under independence---a result
with no series counterpart.  Yet no paper has derived a closed-form
maximum likelihood estimator or EM algorithm for the component rates
of an exponential parallel system from system-level data.
\citet{SarhanElBassiouny2003} treat the exponential parallel system
under a masking framing, deriving MLE and Bayes estimators
numerically, but do not establish identifiability or exploit the
inclusion-exclusion structure.  \citet{CramerNavarro2015} show that
component failure times from a coherent system can be represented as
progressively Type-II censored order statistics---a result that
specializes to our censoring equivalence for the parallel case---but
do not derive closed-form exponential MLE from it.  Other approaches
are nonparametric and signature-based
\citep{BhattacharyaSamaniego2010, NgNavarroBalakrishnan2017,
YangNgBalakrishnan2019, Asadi2020}, Bayesian without closed-form
posteriors \citep{BheringPereiraPolpo2014}, or focused on
non-estimation aspects of system reliability.
Recent work has extended the problem in complementary directions:
\citet{QuNgMoon2024a,QuNgMoon2024b} develop nonparametric tests for
component distribution homogeneity from system data,
\citet{BiswasGangulyMitra2025} propose semiparametric models for
load-sharing parallel systems, and \citet{Jasinski2023,Jasinski2024}
derive closed-form MLE for discrete (geometric) lifetimes in
$k$-out-of-$n$ systems.  None of these addresses the continuous
exponential case with the censoring perspective developed here.

The key obstacle has been a misidentification of the incomplete-data
structure.  Several authors
\citep{LiuShiZhangBai2017, RodriguesPereiraBheringPolpo2018,
Sarhan2001} have transplanted the \emph{masking} framework
(conditions C1--C3) from series to parallel systems.  But masking
presupposes uncertainty about \emph{which} component caused system
failure---a question that is meaningful for series systems (where one
component fails and the rest survive) but structurally meaningless for
parallel systems (where every component has failed by the time the
system fails).

The natural incomplete-data problem for parallel systems is not masking
but \emph{censoring}.  When the system fails at time $T = t$, we know
that all components have failed and that exactly one failed at time~$t$
(the last survivor), but we do not know \emph{when} the other
components failed---only that they failed before~$t$.  This is
component-level left-censoring with latent identity, and it is the
correct framing for estimation.  The censoring perspective was
introduced by \citet{BheringPereiraPolpo2014} in a Bayesian Weibull
context, but its frequentist consequences have not been developed.

\subsection*{Contributions}

We formalize the censoring equivalence and exploit it for exponential
parallel systems with independent components.

\begin{enumerate}
  \item \textbf{Censoring equivalence} (\Cref{prop:censoring-equiv}).
    We prove that observing $T = \max(T_1, \ldots, T_m)$ is
    informationally equivalent to observing one exact and $m-1$
    left-censored component lifetimes with latent identity $J =
    \argmax_j T_j$.

  \item \textbf{Closed-form likelihood via inclusion-exclusion}
    (\Cref{lem:ie,lem:density-decomp,lem:integral-wj}).
    The system density decomposes as $f_T(t) = \sum_j w_j(t)$, where
    each component weight $w_j$ and its integrals have closed forms
    through the inclusion-exclusion expansion of the system CDF.
    This yields a tractable log-likelihood for exact, right-censored,
    and interval-censored system observations without numerical
    quadrature.

  \item \textbf{Elementary identifiability proof}
    (\Cref{thm:identifiability,cor:duality}).
    We provide a self-contained proof that the distribution of
    $T = \max T_j$ uniquely determines the component rates, by
    induction on the number of components via reversed hazard rate
    tail extraction.  A duality corollary formalizes the contrast:
    series systems determine only $\sum \lambda_j$; parallel systems
    determine all individual~$\lambda_j$.

  \item \textbf{EM algorithm with closed-form steps}
    (\Cref{prop:em-algorithm}).
    The censoring equivalence yields an EM algorithm whose E-step
    uses the posterior $P(J = j \mid T = t)$ and the truncated
    exponential mean, and whose M-step is a simple reciprocal.
    Both steps are in closed form---no simulation or numerical
    integration is needed.

  \item \textbf{Monte Carlo validation and information efficiency}
    (\Cref{sec:simulations}).
    Simulations confirm consistency and nominal coverage, and reveal
    an information asymmetry governed by $P(T_j = T)$: slow components
    are estimated precisely, while fast components carry little
    system-level information.  Fisher information analysis
    quantifies this: the per-component efficiency is approximately
    proportional to $P(T_j = T)$, and the joint efficiency
    collapses exponentially with~$m$.

  \item \textbf{$k$-out-of-$n$ generalization}
    (\Cref{prop:kofn-censoring}).
    The censoring equivalence extends to $k$-out-of-$n$ systems:
    observing $T = T_{(n-k+1)}$ yields one exact, $n-k$ left-censored,
    and $k-1$ right-censored component observations with latent identity.
    This unifies the parallel ($k=1$) and series ($k=n$) cases and
    reveals a monotone information structure as~$k$ varies.
\end{enumerate}

\subsection*{Outline}

\Cref{sec:model} introduces the model and the censoring equivalence.
\Cref{sec:likelihood} develops the inclusion-exclusion machinery and
establishes identifiability.
\Cref{sec:estimation} presents the direct MLE and EM algorithm.
\Cref{sec:simulations} reports simulation results, including a Fisher
information comparison.
\Cref{sec:discussion} extends the censoring equivalence to
$k$-out-of-$n$ systems and discusses limitations and future directions.
